% filepath: /Users/rodrigo/dev/uc3m/template_abnt.tex
\documentclass[a4paper,12pt]{article}
\usepackage[brazil]{babel}
\usepackage[utf8]{inputenc}
\usepackage[T1]{fontenc}
\usepackage{amsmath,amssymb}
\usepackage{graphicx}
\usepackage{indentfirst}
\usepackage{cite}
\usepackage{hyperref}
\usepackage{geometry}
\geometry{left=3cm,right=2cm,top=3cm,bottom=2cm}

% Configurações de parágrafo
\setlength{\parindent}{1.25cm}
\setlength{\parskip}{0.2cm}

% Configurações de seção
\usepackage{titlesec}
\titleformat{\section}{\normalfont\bfseries\uppercase}{\thesection}{1em}{}
\titleformat{\subsection}{\normalfont\bfseries}{\thesubsection}{1em}{}
\titleformat{\subsubsection}{\normalfont\bfseries}{\thesubsubsection}{1em}{}

% Capa
\title{Título do Trabalho}
\author{Nome do Autor}
\date{Data}

\begin{document}

\maketitle
\thispagestyle{empty} % Remove numeração da capa

\newpage
\begin{center}
    \textbf{Resumo}
\end{center}
\noindent O resumo deve conter entre 150 e 500 palavras, apresentando de forma concisa os objetivos, metodologia, resultados e conclusões do trabalho.

\begin{center}
    \textbf{Palavras-chave}: palavra1, palavra2, palavra3.
\end{center}

\newpage

\tableofcontents

\newpage

\section{Introdução}
A introdução deve apresentar o tema do trabalho, os objetivos e a justificativa.

\section{Revisão da Literatura}
A revisão da literatura deve apresentar os principais trabalhos já realizados sobre o tema.

\section{Metodologia}
Descreva a metodologia utilizada para a realização do trabalho.

\section{Resultados}
Apresente os resultados obtidos com a pesquisa.

\section{Discussão}
Discuta os resultados obtidos, comparando-os com a literatura existente.

\section{Conclusão}
Apresente as conclusões do trabalho, destacando as contribuições e sugestões para trabalhos futuros.

\newpage

\begin{thebibliography}{99}
\bibitem{exemplo1} SOBRENOME, Nome. \textit{Título do Livro}. Edição. Local: Editora, Ano.
\bibitem{exemplo2} SOBRENOME, Nome. Título do artigo. \textit{Nome da Revista}, Local, v. volume, n. número, p. página inicial-final, mês ano.
\end{thebibliography}

\end{document}